\documentclass[12pt]{article}

\usepackage[a4paper, total={6in, 8in}]{geometry}
\usepackage{textcomp}

\begin{document}
\setlength{\parindent}{0pt}

\def \t {\textrightarrow}
\def \v {\vspace{1em}}
\def \bi {\begin{itemize}}
\def \ei {\end{itemize}}
\def \s[#1] {\section*{#1}}
\def \ss[#1] {\subsection*{#1}}
\def \sss[#1] {\subsubsection*{#1}}

\s[Appunti 21 marzo su foglio a penna]

\s[Architettura fascista]
\ss[Casa del fascio - Terragni]
Capolavoro del razionalismo fascista, a Como \\
Terragni è lo stesso della casa Rustici a Milano \t studia a como ma si laurea al politecnico di milano \\
Sara sulla nave della Carta d'Atene \t scrive riviste, e lavora soprattutto a Como (questa opera, poi fa Novum Covum, asilo di Sant'Elia) \\
Poi va in guerra (seconda guerra mondiale) \t andrà anche in Russia \t muore giovane nel 43 anni per trombosi cerebrale \\
Diventa la sede del partito fascista \t ora è sede della guardia di finanza \\
Dal punto di vista funzionale ci sono quindi gli uffici e la sala delle adunanze

\v

Al piano terra ci sono corridoi, a primo piano ballatoi che si affacciano sulla sala delle adunanze \t e ballatoi collegano i diversi uffici \\
Non c'è un trattamento monumentale di una facciata rispetto alle altre \t no facciata principale vera \\
4 piani, che si organizza intorno a una corte decentrata \t la corte è la sala delle adunanza \\
Collegamento verticale è una scala \\
Struttura a corte richiama alla tradizione architettonica italiana \t tipica del palazzo rinascimentale \\
La progettazione si basa sulla geometria \t usa il quadrato come forma base \t rimanda a rinascimento con Alberti \t e il palazzo è un mezzo cubo

\ss[Struttura]
Punto di vista strutturale ci sono pilastri e travi in calcestruzzo armato \t e sono a vista, sono visibili + tamponamenti in laterizio \\
La copertura della sala delle adunanza è il vetro cemento \t brevettato a inizio 900, usato principalmente in edifici industriali \\
Il vetro cemento \t sono due formelle di vetro, con intercapedine di aria \t e poi vengono unite con cemento \t proprietà è che fa passare la luce, ma impedisce di vedere bene all'interno \t ed è molto resistente \\
Materiale nuovo che usa \\
Il rivestimento esterno è tutto in marmo bianco \t il movimento moderno non usa materiali preziosi come il marmo \t anche in una struttura cosi geometrica ed essenziale da + monumentalita, ma lui giustifica dicendo che marmo è molto resistente \

\ss[Interno]
L'interno è intonacato in verde pastello \t ora invece è bianca \\
La scala è molto essenziale, il sistema strutturale è esibito \t nessun tipo di decorazione, ne interno ne esterno

\v

I prospetti sono tutti diversi \t no facciata principale \t sono tutti con marmo e studiati su base geometrica \\
Suddivide geometricamente dei rettangoli, alti 4 e lunghi 7 \t lascia i due rettangoli laterali coperti, mentre negli altri ci sono logge con struttua portante a vista \\
Progettazione su base geometrica evidente
